\chapter{Statistiek en kansrekening}\label{ch:g9:statproba}

Gegevens en toeval zijn overal: sportuitslagen, weersvoorspellingen, spelletjes.
Dit hoofdstuk leert hoe je een reeks getallen samenvat met haar \omterm{def:g9:statproba:indicators}{gemiddelde}, haar
\omterm{def:g9:statproba:indicators}{mediaan} en haar \omterm{def:g9:statproba:indicators}{variatiebreedte}, en hoe je de \omterm{def:g9:statproba:probability}{kans} van eenvoudige kansproeven
berekent --- ook van proeven in twee stappen, die je met boomdiagrammen aanpakt.
Beide onderwerpen worden veel verder uitgewerkt in het bovenbouwvolume van deze
reeks.

\section{Gegevens samenvatten}

\begin{definition}[Gemiddelde, mediaan, variatiebreedte]\label{def:g9:statproba:indicators}
Voor een reeks van $N$ getallen:
\begin{itemize}
  \item is het \emph{gemiddelde}\index{gemiddelde} de \omterm{def:g1:addition:def}{som} van alle waarden gedeeld
        door $N$;
  \item is de \emph{mediaan}\index{mediaan} een waarde die de gesorteerde reeks in
        twee \omterm{def:g3:division:half}{helften} van gelijke grootte verdeelt: voor \omterm{def:g3:numbers:evenodd}{oneven} $N$ de middelste
        waarde, voor \omterm{def:g3:numbers:evenodd}{even} $N$ het midden van de twee middelste waarden;
  \item is de \emph{variatiebreedte}\index{variatiebreedte} de grootste waarde min
        de kleinste.
\end{itemize}
\end{definition}

\begin{example}\label{ex:g9:statproba:indicators}
De cijfers van een leerling: $8,\ 12,\ 9,\ 15,\ 11$.
\begin{enumerate}
  \item \omterm{def:g9:statproba:indicators}{Gemiddelde}: $\dfrac{8 + 12 + 9 + 15 + 11}{5} = \dfrac{55}{5} = 11$.
  \item \omterm{def:g9:statproba:indicators}{Mediaan}: sorteer de reeks: $8, 9, 11, 12, 15$; de middelste (3de) waarde is
        $11$.
  \item \omterm{def:g9:statproba:indicators}{Variatiebreedte}: $15 - 8 = 7$.
\end{enumerate}
Met een zesde cijfer van $17$: gesorteerd $8, 9, 11, 12, 15, 17$, wordt de \omterm{def:g9:statproba:indicators}{mediaan}
het midden van $11$ en $12$, dat is $11.5$, en wordt het \omterm{def:g9:statproba:indicators}{gemiddelde}
$\frac{72}{6} = 12$.
\end{example}

\begin{example}[Gewogen gemiddelde]\label{ex:g9:statproba:weighted}
De schoenmaten die op één dag verkocht zijn:
\begin{center}
\renewcommand{\arraystretch}{1.2}
\begin{tabular}{c|ccccc}
maat & $38$ & $39$ & $40$ & $41$ & $42$ \\
\hline
aantal & $3$ & $7$ & $6$ & $3$ & $1$ \\
\end{tabular}
\end{center}
De \omterm{def:g9:statproba:indicators}{gemiddelde} maat is
\[
\frac{3 \times 38 + 7 \times 39 + 6 \times 40 + 3 \times 41 + 1 \times 42}
{3 + 7 + 6 + 3 + 1}
= \frac{114 + 273 + 240 + 123 + 42}{20}
= \frac{792}{20} = 39.6 .
\]
\end{example}

\begin{remark}
Het \omterm{def:g9:statproba:indicators}{gemiddelde} en de \omterm{def:g9:statproba:indicators}{mediaan} kunnen sterk verschillen. In de reeks
$10, 10, 10, 10, 60$ is het \omterm{def:g9:statproba:indicators}{gemiddelde} $20$ (omhooggetrokken door de waarde $60$),
terwijl de \omterm{def:g9:statproba:indicators}{mediaan} $10$ is. Vraag je altijd af welke maat de situatie beter
weergeeft.
\end{remark}

\section{Kansrekening}

\begin{definition}[Kans]\label{def:g9:statproba:probability}
Een kansproef heeft verschillende mogelijke \emph{uitkomsten}; een
\emph{gebeurtenis} is een verzameling uitkomsten. Elke uitkomst krijgt een
\emph{kans}\index{kans}: een getal tussen $0$ en $1$, waarbij alle kansen van de
uitkomsten samen $1$ zijn; de kans $\P(A)$ van een gebeurtenis $A$ is de \omterm{def:g1:addition:def}{som} van de
kansen van haar uitkomsten. Zijn alle $n$ uitkomsten even waarschijnlijk, dan is
\[
\P(A) = \frac{\text{aantal gunstige uitkomsten}}{n}.
\]
\end{definition}

\begin{example}\label{ex:g9:statproba:spinner}
Een rad is in $8$ gelijke sectoren verdeeld: $3$ rode, $3$ blauwe, $2$ groene. Elke
sector heeft \omterm{def:g9:statproba:probability}{kans} $\frac18$, dus is
\[
\P(\text{rood}) = \frac38, \qquad
\P(\text{groen}) = \frac28 = \frac14, \qquad
\P(\text{niet groen}) = 1 - \frac14 = \frac34 .
\]
\end{example}

\begin{proposition}[Complement]\label{prop:g9:statproba:complement}
Voor elke gebeurtenis $A$ voldoet de \emph{complementaire gebeurtenis} $\bar A$
(``$A$ treedt niet op'') aan
\[
\P(\bar A) = 1 - \P(A).
\]
\end{proposition}

\begin{proof}
Elke uitkomst zit in precies een van $A$ en $\bar A$, dus telt
$\P(A) + \P(\bar A)$ de kansen van alle uitkomsten op, en dat is $1$.
\end{proof}

\section{Proeven in twee stappen}

\begin{method}[Boomdiagrammen]\label{met:g9:statproba:tree}
Voor een proef in twee stappen:
\begin{enumerate}
  \item teken één tak per uitkomst van de eerste stap, en zet elke tak voort met de
        uitkomsten van de tweede stap; schrijf elke \omterm{def:g9:statproba:probability}{kans} bij haar tak;
  \item vermenigvuldig de kansen langs een pad om de \omterm{def:g9:statproba:probability}{kans} van dat pad te krijgen;
  \item tel de kansen op van alle paden die een gebeurtenis vormen.
\end{enumerate}
\end{method}

\begin{example}\label{ex:g9:statproba:tree}
Een zak bevat $2$ rode en $3$ zwarte schijfjes. Trek één schijfje, noteer zijn
kleur, \emph{leg het terug} en trek opnieuw. Elke trekking geeft rood met \omterm{def:g9:statproba:probability}{kans}
$\frac25$ en zwart met \omterm{def:g9:statproba:probability}{kans} $\frac35$.

\begin{omfigure}
\begin{tikzpicture}[scale=0.95]
  \coordinate (root) at (0,0);
  \coordinate (R) at (2.2,1.0);
  \coordinate (B) at (2.2,-1.0);
  \coordinate (RR) at (4.7,1.55);
  \coordinate (RB) at (4.7,0.45);
  \coordinate (BR) at (4.7,-0.45);
  \coordinate (BB) at (4.7,-1.55);
  \draw (root) -- (R) node[midway, above, font=\small, sloped] {$\frac25$};
  \draw (root) -- (B) node[midway, below, font=\small, sloped] {$\frac35$};
  \draw (R) -- (RR) node[midway, above, font=\small, sloped] {$\frac25$};
  \draw (R) -- (RB) node[midway, below, font=\small, sloped] {$\frac35$};
  \draw (B) -- (BR) node[midway, above, font=\small, sloped] {$\frac25$};
  \draw (B) -- (BB) node[midway, below, font=\small, sloped] {$\frac35$};
  \node[omThm, font=\small, above left=-2pt] at (R) {R};
  \node[font=\small, below left=-2pt] at (B) {Z};
  \node[right, font=\small] at (RR)
    {RR:\ $\frac25 \times \frac25 = \frac{4}{25}$};
  \node[right, font=\small] at (RB)
    {RZ:\ $\frac25 \times \frac35 = \frac{6}{25}$};
  \node[right, font=\small] at (BR)
    {ZR:\ $\frac35 \times \frac25 = \frac{6}{25}$};
  \node[right, font=\small] at (BB)
    {ZZ:\ $\frac35 \times \frac35 = \frac{9}{25}$};
\end{tikzpicture}

{\small De boom van de twee trekkingen met teruglegging: vermenigvuldig langs de
takken, en ga na dat de vier paden samen $1$ geven.}
\end{omfigure}

\omterm{def:g9:statproba:probability}{Kans} op twee schijfjes van dezelfde kleur:
$\frac{4}{25} + \frac{9}{25} = \frac{13}{25}$. \omterm{def:g9:statproba:probability}{Kans} op minstens één rood:
$1 - \P(\text{ZZ}) = 1 - \frac{9}{25} = \frac{16}{25}$.
\end{example}

\begin{remark}[Frequenties naderen kansen]
Een eerlijke dobbelsteen $6000$ keer gooien geeft \emph{ongeveer} $1000$ zessen ---
niet precies. Naarmate het aantal herhalingen groeit, komen de waargenomen
\omterm{def:g7:stats:frequency}{frequenties} steeds dichter bij de kansen: dat is wat de kansrekening het juiste
gereedschap maakt om herhaalde proeven te beschrijven (een verhaal dat in het
bovenbouwvolume van deze reeks wordt uitgewerkt).
\end{remark}

\section{Oefeningen}

\begin{exercise}[$\star$]\label{exo:g9:statproba:1}
Bereken het \omterm{def:g9:statproba:indicators}{gemiddelde}, de \omterm{def:g9:statproba:indicators}{mediaan} en de \omterm{def:g9:statproba:indicators}{variatiebreedte} van de reeks
$14,\ 9,\ 17,\ 9,\ 12,\ 11,\ 12$.
\end{exercise}

\begin{exercise}[$\star$]\label{exo:g9:statproba:2}
De temperaturen op de middag waren gedurende een week $19,\ 21,\ 24,\ 18,\ 22,\
25,\ 19$ (graden). Bereken het \omterm{def:g9:statproba:indicators}{gemiddelde} (op een tiende) en de \omterm{def:g9:statproba:indicators}{mediaan}.
\end{exercise}

\begin{exercise}[$\star$]\label{exo:g9:statproba:3}
Het aantal doelpunten dat een ploeg in $10$ wedstrijden maakte:
\begin{center}
\renewcommand{\arraystretch}{1.2}
\begin{tabular}{c|cccc}
doelpunten & $0$ & $1$ & $2$ & $3$ \\
\hline
wedstrijden & $2$ & $4$ & $3$ & $1$ \\
\end{tabular}
\end{center}
Bereken het \omterm{def:g9:statproba:indicators}{gemiddelde} aantal doelpunten per wedstrijd, en de \omterm{def:g9:statproba:indicators}{mediaan}.
\end{exercise}

\begin{exercise}[$\star$]\label{exo:g9:statproba:4}
Een eerlijke dobbelsteen wordt één keer gegooid. Bereken de \omterm{def:g9:statproba:probability}{kans} op: ``een $3$
gooien''; ``een \omterm{def:g3:numbers:evenodd}{even getal} gooien''; ``minstens een $3$ gooien''; ``geen $6$
gooien''.
\end{exercise}

\begin{exercise}[$\star$]\label{exo:g9:statproba:5}
Een doos bevat $12$ ballen: $5$ witte, $4$ rode en $3$ groene; er wordt één bal
willekeurig getrokken. Bereken $\P(\text{wit})$, $\P(\text{rood of groen})$ en
$\P(\text{niet wit})$. Wat valt je op aan de laatste twee?
\end{exercise}

\begin{exercise}[$\star\star$]\label{exo:g9:statproba:6}
Er wordt een munt opgegooid en daarna een eerlijke dobbelsteen gegooid. Teken de
boom (of tel de uitkomsten) en bereken de \omterm{def:g9:statproba:probability}{kans} op kop \emph{en} een $6$; en daarna
op kop \emph{of} een $6$ (of beide).
\end{exercise}

\begin{exercise}[$\star\star$]\label{exo:g9:statproba:7}
In de zak van \cref{ex:g9:statproba:tree} ($2$ rode, $3$ zwarte) worden de twee
trekkingen nu \emph{zonder} teruglegging gedaan. Teken de nieuwe boom (let op: de
kansen op het tweede niveau veranderen) en bereken de \omterm{def:g9:statproba:probability}{kans} op twee zwarte
schijfjes, en daarna op twee schijfjes van verschillende kleur.
\end{exercise}

\begin{exercise}[$\star\star$]\label{exo:g9:statproba:8}
Na $9$ wedstrijden haalt een basketbalspeelster een \omterm{def:g9:statproba:indicators}{gemiddelde} van $14$ punten per
wedstrijd.
\begin{enumerate}
  \item Hoeveel punten heeft ze in totaal gemaakt?
  \item Hoeveel punten moet ze in de $10$de wedstrijd maken om op een \omterm{def:g9:statproba:indicators}{gemiddelde}
        van $15$ te komen?
\end{enumerate}
\end{exercise}

\begin{exercise}[$\star\star\star$]\label{exo:g9:statproba:9}
Een spel: gooi met twee eerlijke dobbelstenen; je wint als de twee uitkomsten
gelijk zijn (``een dubbel'').
\begin{enumerate}
  \item Bereken de \omterm{def:g9:statproba:probability}{kans} om te winnen.
  \item Je speelt twee keer (onafhankelijke spelletjes). Bereken met een boom de
        \omterm{def:g9:statproba:probability}{kans} om minstens één keer te winnen.
\end{enumerate}
\end{exercise}

\section{Opgave: de ruïneuze weddenschappen van de Chevalier}

\begin{problem}[{Weekendopgave --- het geschil over gokken uit 1654 dat de
kansrekening deed ontstaan, en het geboortedagtoeval dat elke klas om de tuin
leidt}]
\label{pb:g9:statproba:1}
In 1654 werd een Franse edelman en verstokte gokker, de Chevalier de Méré, rijk met
één weddenschap op dobbelstenen en begon hij te verliezen op een andere, die hij
gelijkwaardig achtte. Verbijsterd vroeg hij het aan de wiskundige Blaise Pascal;
Pascal schreef naar Pierre de Fermat; en hun briefwisseling stichtte de theorie van
de kansrekening. Deze opgave speelt de hele zaak opnieuw met het gereedschap van dit
hoofdstuk --- even waarschijnlijke uitkomsten, bomen en de complementregel
(\cref{prop:g9:statproba:complement}) --- en eindigt met het meest tegenintuïtieve
toeval van allemaal.

\medskip
\noindent\textbf{Deel I --- Dobbelstenen, eerlijk geteld.}
\begin{enumerate}
  \item Gooi twee keer met een eerlijke dobbelsteen. Bereken met een boom (of een
        tabel van $6 \times 6$) de \omterm{def:g9:statproba:probability}{kans} op \emph{geen} zes in de twee worpen, en
        leid de \omterm{def:g9:statproba:probability}{kans} op minstens één zes af.
  \item Een klasgenoot beweert: ``één \omterm{def:g9:statproba:probability}{kans} op zes per worp, twee keer, dus
        $\frac16 + \frac16 = \frac13$.'' Vergelijk met vraag 1 en wijs precies de
        uitkomsten aan die zijn optelling twee keer meetelt.
  \item Met twee dobbelstenen gooien en optellen: bereken de \omterm{def:g9:statproba:probability}{kans} op een \omterm{def:g1:addition:def}{som} van
        $7$ en op een \omterm{def:g1:addition:def}{som} van $2$. Waarom is $7$ de favoriet van de gokkers?
  \item Een oud geschil: is met twee dobbelstenen een \omterm{def:g1:addition:def}{som} van $9$ of een \omterm{def:g1:addition:def}{som} van
        $10$ waarschijnlijker? Tel de uitkomsten en beslecht het.
  \item Veralgemeen vraag 1: leg uit waarom de \omterm{def:g9:statproba:probability}{kans} op minstens één zes in $n$
        worpen
        \[
        1 - \left(\frac56\right)^{\!n}
        \]
        is, doordat de takken zonder zes zich van niveau tot niveau
        vermenigvuldigen.
\end{enumerate}

\medskip
\noindent\textbf{Deel II --- De twee weddenschappen van de Chevalier.}
\begin{enumerate}[resume]
  \item Eerste weddenschap, tegen gelijke inzet: \emph{minstens één zes in vier
        worpen met een dobbelsteen}. Bereken $\left(\frac56\right)^4$ als \omterm{def:g9:fractions:fraction}{breuk} en
        als \omterm{def:g6:decimals:places}{decimaal getal}, en de \omterm{def:g9:statproba:probability}{kans} dat de Chevalier wint. Was de weddenschap
        goed?
  \item De Chevalier redeneerde zelf: ``vier worpen van elk $\frac16$: $\frac46$
        \omterm{def:g9:statproba:probability}{kans}.'' Hier gaf dat bijna het juiste antwoord --- maar drijf het door naar
        $7$ worpen: welke ongerijmdheid levert het op, en welke fout uit vraag 2
        herhaalt het?
  \item Tweede weddenschap, die hij op grond van \omterm{def:g9:linfunc:linear}{evenredigheid} gelijkwaardig achtte
        (``$24$ worpen van elk $\frac1{36}$: weer $\frac{24}{36} = \frac46$''):
        \emph{minstens één dubbele zes in $24$ worpen met twee dobbelstenen}.
        Bereken de werkelijke \omterm{def:g9:statproba:probability}{kans} om te winnen,
        $1 - \left(\frac{35}{36}\right)^{24}$ (met de rekenmachine). Waarom werd de
        Chevalier langzaam geruïneerd?
  \item Hoeveel worpen met twee dobbelstenen zouden zijn voordeel hebben
        hersteld? Toets $n = 25$ met de rekenmachine en besluit.
  \item Formuleer de les in twee zinnen: wat is er in het algemeen verkeerd aan een
        \omterm{def:g9:statproba:probability}{kans} met het aantal pogingen vermenigvuldigen --- en welk juist gereedschap
        vervangt die verleiding?
\end{enumerate}

\medskip
\noindent\textbf{Deel III --- \omterm{def:g9:statproba:indicators}{Gemiddelden}, medianen, geboortedagen.}
\begin{enumerate}[resume]
  \item Een klein bedrijf betaalt negen werknemers $2\,000$ euro per maand en de
        directeur $20\,000$. Bereken het \omterm{def:g9:statproba:indicators}{gemiddelde} en de \omterm{def:g9:statproba:indicators}{mediaan} van de lonen
        (\cref{def:g9:statproba:indicators}). Welk getal zou de vacature eerlijk
        moeten vermelden?
  \item Stel een lijst van vijf toetscijfers op met \omterm{def:g9:statproba:indicators}{gemiddelde} $12$ en \omterm{def:g9:statproba:indicators}{mediaan}
        $15$. Wat zegt dat paar (\omterm{def:g9:statproba:indicators}{gemiddelde} onder de \omterm{def:g9:statproba:indicators}{mediaan}) over de vorm van de
        cijfers?
  \item Het geboortedagprobleem: wat is in een klas van $23$ leerlingen de \omterm{def:g9:statproba:probability}{kans} dat
        minstens twee dezelfde geboortedag hebben? Stel het complement op (alle
        $23$ geboortedagen verschillend):
        \[
        \frac{365}{365} \times \frac{364}{365} \times
        \frac{363}{365} \times \dots \times \frac{343}{365} .
        \]
        Bereken het \omterm{def:g2:mult:def}{product} van de eerste drie factoren, beschrijf hoe de factoren
        evolueren, en antwoord op de vraag, gegeven dat het volledige \omterm{def:g2:mult:def}{product}
        ongeveer $0.493$ is. De meeste mensen verwachten ``heel onwaarschijnlijk'':
        wat zegt de wiskunde?
  \item De intuïtie hersteld: hoeveel \emph{paren} leerlingen kunnen in een klas
        van $23$ dezelfde geboortedag hebben (het tellen van de handdrukken uit
        \cref{pb:g6:wholes:1})? Leg in één zin uit waarom dat getal, en niet $23$
        zelf, de verrassing veroorzaakt.
  \item Slot: beschrijf een proef om een van beide resultaten na te gaan --- de
        weddenschap van de Méré met een echte dobbelsteen, of het
        geboortedagprobleem over verschillende klassen --- en formuleer zorgvuldig
        wat je van de waargenomen \omterm{def:g7:stats:frequency}{frequenties} verwacht naarmate het aantal
        herhalingen groeit. (Die verwachting heeft een naam, de wet van de grote
        aantallen; haar bewijs wacht in het bovenbouwvolume.)
\end{enumerate}
\end{problem}
